% Replaces Section 6 "Theory" of the draft (PAPER_TODO 1.2-1.6).  Source of truth: docs/paper/THEORY.md, Sections 1-5.
% Numerical checks: tools/theory_checks.py (workflow "Gatekeeper theory checks").
\section{Theory}
\label{sec:theory}

We claim no new statistics: the guarantee is assembled from betting e-processes~\cite{waudbysmith2024betting},
Ville's inequality~\cite{ville1939} and a union bound.  What is new is the estimand, the per-cell ledger that covers
compositions and re-admission, the proof that the shift term is unavoidable, and the exact characterisation of composition
risk.

\subsection{Per-cell evidence}

\begin{lemma}[Betting e-process]\label{lem:eprocess}
Fix a cell $c$ and let $Y_i=-Z_i$.  Under $H_0^c:\ \mathbb{E}[Y]\le\theta$ (the cell is not admissible), the process
$K_0=1$, $K_n=\prod_{i\le n}\bigl(1+\lambda_i(Y_i-\theta)\bigr)$ with predictable $\lambda_i\in[0,1/(1+\theta))$ satisfies
$\Pr[\exists n:\ K_n\ge 1/\beta]\le\beta$ for every $\beta\in(0,1)$.
\end{lemma}

\begin{proof}
Since $Y_i\ge-1$, every factor is at least $1-\lambda_i(1+\theta)>0$.  By Assumption~\ref{ass:trials} and predictability,
$\mathbb{E}[K_i\mid\mathcal{F}_{i-1}]=K_{i-1}\bigl(1+\lambda_i(\mathbb{E}[Y_i]-\theta)\bigr)\le K_{i-1}$, so $(K_n)$ is a
nonnegative supermartingale with $K_0=1$, and Ville's inequality gives the claim.
\end{proof}

The bound holds at any data-dependent stopping time, including the end of the validation budget, which is why a fixed-sample
interval would not do.

\subsection{Library-wide admission guarantee}

\begin{theorem}[Admission over a declared scope family]\label{thm:admission}
Let $\mathcal{C}$ be the set of cells, compositions included, and let a ledger $\alpha_{c,k}\ge0$ be fixed in advance for
every cell $c$ and epoch $k\ge0$ (epoch $k\ge1$ is the $k$-th re-admission after retirement), with
$\sum_{c\in\mathcal{C}}\sum_{k\ge0}\alpha_{c,k}\le\alpha$; for instance $\alpha_{c,k}=\alpha_c\cdot 6/(\pi^2(k+1)^2)$ with
$\sum_c\alpha_c\le\alpha$.  Suppose each epoch uses fresh tasks and admits $c$ only when its e-process
(Lemma~\ref{lem:eprocess}) reaches $1/\alpha_{c,k}$, while the choice of cells, their order and the stopping times may depend
on all data.  Under Assumptions~\ref{ass:trials} and~\ref{ass:shift},
\[
  \Pr\bigl[\exists\text{ admitted }(c,k):\ \mathrm{STR}_{Q}(c)>-\theta+2\xi_{s(c)}\bigr]\le\alpha .
\]
In particular, if $2\xi_s<\theta$ for every scope, the probability of admitting any harmful cell is at most $\alpha$.
\end{theorem}

\begin{proof}
For every $(c,k)$ let $A_{c,k}=\{\sup_n K^{c,k}_n\ge1/\alpha_{c,k}\}$, defined on the whole infinite task sequence of that
epoch.  An admission of $(c,k)$ implies $A_{c,k}$; adaptive selection and stopping can only forgo admissions.  Hence the
event that some cell with $H_0^c$ true is admitted is contained in $\bigcup_{(c,k):H_0^c}A_{c,k}$, whose probability is at
most $\sum\alpha_{c,k}\le\alpha$ by Lemma~\ref{lem:eprocess} and a union bound, which needs no independence across cells.
Finally, by Lemma~\ref{lem:shift}, $\mathrm{STR}_{Q}(c)>-\theta+2\xi_s$ implies $\mathrm{STR}_{P}(c)>-\theta$, that is,
$H_0^c$.
\end{proof}

\begin{lemma}[Shift]\label{lem:shift}
If $f$ takes values in an interval of length $b$, then $|\mathbb{E}_Qf-\mathbb{E}_Pf|\le b\cdot\mathrm{TV}(P,Q)$.  Since
$Z\in[-1,1]$, $|\mathrm{STR}_Q(c)-\mathrm{STR}_P(c)|\le2\,\mathrm{TV}(P_s,Q_s)\le2\xi_s$, and the constant $2$ is attained.
\end{lemma}

\begin{proof}
Centre $f$ so that $|g|\le b/2$; then $|\mathbb{E}_Qg-\mathbb{E}_Pg|\le(b/2)\sum_x|Q(x)-P(x)|=b\,\mathrm{TV}(P,Q)$.  For
attainment, take task types $u$ with $Z=1$ and $d$ with $Z=-1$, $P(u)=0.1$, $Q(u)=0.2$: then $\mathrm{TV}=0.1$ and the risks
differ by $0.2$.
\end{proof}

\paragraph{Why the budget is per cell.}
Allocating one $\alpha_s$ per \emph{scope} and letting every skill in the scope use the threshold $1/\alpha_s$ does not
control the family-wise error: with ten skills at the null boundary, $\theta=\alpha=0.1$ and $300$ trials, the family-wise
false admission rate of our e-process is $0.45$, against $0.05$ with the per-cell budget $\alpha/10$
(Section~\ref{sec:checks}).

\subsection{Scope and shift are necessary}

\begin{assumption}[Rich target class]\label{ass:rich}
For every value of the features a procedure observes (fingerprints, predicates, trial outcomes on the observed
repositories), there is a repository with those features whose task losses are arbitrary in $[0,1]$.
\end{assumption}

\begin{proposition}[No guarantee without target evidence or bounded shift]\label{prop:nofree}
\textup{(a)} Let a procedure decide from data of observed repositories $\mathcal{R}_{\mathrm{obs}}$ only, and admit a skill
$S$ with positive probability.  Under Assumption~\ref{ass:rich} there is a repository $R^\dagger\notin\mathcal{R}_{\mathrm{obs}}$
with $\mathrm{STR}(S,\{R^\dagger\})=1$ on which the admission probability is unchanged.
\textup{(b)} For every $\theta\in(0,1)$, $\delta>0$ with $\theta+\delta<1$, $\xi\in(0,1/2]$ and every procedure, there are
$P,Q$ with $\mathrm{TV}(P,Q)=\xi$ and a cell $c$ with $\mathrm{STR}_P(c)=-\theta-\delta$ and
$\mathrm{STR}_Q(c)=\mathrm{STR}_P(c)+2\xi$, whose admission probability does not depend on $Q$.
\end{proposition}

\begin{proof}
(a) Take $R^\dagger$ that agrees with all observed data and has $\ell(x,\pi)=1$, $\ell(x,\varnothing)=0$ for every task.
(b) Take task types $u$ ($Z=1$) and $d$ ($Z=-1$) with $P(d)=(1+\theta+\delta)/2\ge\xi$, and obtain $Q$ by moving mass $\xi$ from
$d$ to $u$.  The procedure only reads tasks drawn from $P$.
\end{proof}

Part (b) shows that no procedure that learns only from $P$ can certify $Q$-risk better than $\mathrm{STR}_P+2\xi$, so the shift
term of Theorem~\ref{thm:admission} cannot be removed or reduced; this is why the guarantee is relative to a declared scope
family and why coarse scopes make it weak.

\subsection{Composition}

\begin{theorem}[Composition risk is not controlled by individual risks]\label{thm:composition}
For $a,b\in[-1,0)$,
\[
  \sup\bigl\{\mathrm{STR}(S_1\circ S_2):\ \mathrm{STR}(S_1)=a,\ \mathrm{STR}(S_2)=b\bigr\}=1+\min(a,b),
\]
and the supremum is attained.  In particular, two skills that each meet the margin $-\theta$ can compose to excess risk
$1-\theta>0$, and no bound on the composite risk improves on the trivial one implied by $\ell\in[0,1]$.
\end{theorem}

\begin{proof}
Upper bound: $\mathrm{STR}(S_1\circ S_2)\le1-\mathbb{E}\ell(\varnothing)$, and $\mathbb{E}\ell(\varnothing)\ge-a$ because
$\mathbb{E}\ell(\pi_1)=\mathbb{E}\ell(\varnothing)+a\ge0$; likewise $\ge-b$.  Attainment: constant losses
$\ell(\varnothing)=-\min(a,b)$, $\ell(\pi_1)=\ell(\varnothing)+a$, $\ell(\pi_2)=\ell(\varnothing)+b$,
$\ell(\pi_1\circ\pi_2)=1$.
\end{proof}

Theorem~\ref{thm:composition} says nothing about footprints.  Screening only pairs whose footprints overlap is therefore sound
only under an additional assumption, which we state and test rather than presume.

\begin{assumption}[Non-interaction of disjoint footprints]\label{ass:noninteract}
If $\mathcal{K}\cap\mathcal{K}'=\varnothing$ then
$\ell(x,\pi\circ\pi')-\ell(x,\varnothing)=[\ell(x,\pi)-\ell(x,\varnothing)]+[\ell(x,\pi')-\ell(x,\varnothing)]$ for all $x$.
\end{assumption}

\begin{corollary}\label{cor:screening}
Under Assumption~\ref{ass:noninteract}, two admissible skills with disjoint footprints compose to
$\mathrm{STR}\le-2\theta$, so only footprint-overlapping pairs need a composite cell; with mean overlap degree $\bar d$ there
are $|\mathcal{L}_s|\,\bar d/2$ such pairs per scope.
\end{corollary}

\subsection{The value of pairing}

\begin{proposition}[Pairing]\label{prop:pairing}
Let $V_\pi=\mathrm{Var}\,\ell(x,\pi)$, $V_\varnothing=\mathrm{Var}\,\ell(x,\varnothing)$ and
$C=\mathrm{Cov}(\ell(x,\pi),\ell(x,\varnothing))$ on the same task.  Estimating $\mathrm{STR}$ from $n$ paired tasks has
variance $(V_\pi+V_\varnothing-2C)/n$; from two independent arms of $n$ tasks each it has variance $(V_\pi+V_\varnothing)/n$.
Both designs cost $2n$ runs, so pairing reduces the number of tasks needed for a given precision by
$R=(V_\pi+V_\varnothing)/(V_\pi+V_\varnothing-2C)$.  Under the additive model $\ell(\varnothing)=T$, $\ell(\pi)=T+S$ with
independent $T,S$, $R=1+2\sigma_T^2/\sigma_S^2$.
\end{proposition}

\begin{proof}
Direct computation of the two variances; under the additive model $V_\pi+V_\varnothing=2\sigma_T^2+\sigma_S^2$ and $C=\sigma_T^2$.
\end{proof}

\subsection{Numerical checks}
\label{sec:checks}
Every statement above is checked numerically by a script in the artifact, using the e-process of our implementation:
per-cell validity, the failure of a per-scope budget, the attained constant $2$ in Lemma~\ref{lem:shift}, the bound of
Theorem~\ref{thm:composition} on random instances, the pairing ratio, and the convergence of the ledger.
